\begin{abstract}
Hardware security mechanisms are burden-allocation instruments: their technical
properties determine who pays for security, when, and how much---as a direct
consequence of design choices that could have been made differently. This reframing
exposes a systematic flaw in standard evaluation practice. Benchmark reporting
measures local overhead precisely but says nothing about where costs land; the
result systematically favors mechanisms that externalize costs over mechanisms
that internalize them, because externalized costs are invisible in overhead figures.

We develop this argument through a synthesis of systems engineering and security
economics---two traditions that analyze the same mechanisms with incompatible
vocabularies and complementary blind spots---and ground it in six case studies
spanning speculative execution controls, memory safety, secure boot, and hardware
random number generation. The product is a semi-formal burden-allocation vocabulary
that makes cost pathways explicit, attributable, and analyzable before deployment.
\end{abstract}